\selectlanguage{ngerman}%
\frenchspacing%
\begin{ledgroupsized}[r]{120mm}%
\footnotesize%
\pstart%
\noindent%
\textbf{Überlieferung:}%
\pend%
\end{ledgroupsized}%
\begin{ledgroupsized}[r]{114mm}%
\footnotesize%
\pstart%
\parindent -6mm%
\makebox[6mm][l]{\textit{L}}%
Notiz:
LH XXXV 14,\;2 Bl.\;1.
Ein Zettel (3\;x\;10 cm),
Ränder beschnitten und zum Teil gerissen;
möglicherweise italienisches Papier (ohne Wasserzeichen).
Zwei vollbeschriebene Seiten.
\pend
\end{ledgroupsized}%
\vspace{4mm}%
\begin{ledgroup}%
\footnotesize%
\pstart%
\noindent%
\textbf{Datierungsgründe}: %
Die Datierung der Notiz N.\;\ref{58216} beruht auf äußeren Eigenschaften des Text\-trägers %
sowie auf inhaltlichen Betrachtungen.
Siehe die ausführliche Begründung auf S.\;\pageref{58216_dat}.%
\pend%
\end{ledgroup}%
\selectlanguage{latin}%
\frenchspacing%
\vspace{8mm}%
\normalsize%
\pstart%
\noindent%
\edtext{%
\bracks{%
\href{\Scan{LH_35_14_02_001r}}%
{1\,r\textsuperscript{o}}%
}
Si duo sint mobilia%
\indexS{mobile e materia uniformi}
aequale spatium occupantia%
\indexS{mobilia occupantia spatium aequale}%
\lbrack,\rbrack\
unumquodque in omnibus suis partibus
materiae uniformis%
\indexS{materia uniformis}
(ut duo aeres pedales%
\indexS{aer pedalis}%
\lbrack)\rbrack;
et ambo facta sint
ex alia quadam materia uniformi prius%
\indexS{materia uniformis}
(ut aere communi)%
\indexS{aer communis}
sed cujus quantitates diversae%
\indexS{quantitas materiae}
intra duo mobilia sint redactae,
(ita ut unus aer pedalis%
\indexS{aer pedalis}
\bracks{%
\href{\Scan{LH_35_14_02_001v}}%
{1\,v\textsuperscript{o}}%
}
duplo plus aeris communis
\lbrack sit\rbrack\
quam alter)
tunc densitas est ea ipsa%
\indexS{densitas aeris}%
\lbrack,\rbrack\
in ipsis mobilibus%
\indexS{densitas in mobili} %
aut aliis conformibus%
\lbrack,\rbrack\
quae sunt
materiae constituentis quantitates%
\indexS{materia constituens}%
\indexS{quantitas materiae}
(seu unus aer pedalis%
\indexS{aer pedalis}
altero dicitur duplo densior%
\indexS{aer densior}%
\indexS{densitas aeris}%
\lbrack)\rbrack.%
}{%
\lemma{\href{\Scan{LH_35_14_02_001r}}{\lbrack1\,r\textsuperscript{o}\rbrack}}%
\Bfootnote{%
\textit{(1)}~Densitas mobilis est ut factum %
\textit{(2)}~Si duo sint mobilia %
\textit{(a)}~aequalis %
\textit{(aa)}~extensionis %
\textit{(bb)}~spatii et materiae uniformis, %
\textit{(b)}~aequale spatium %
\lbrack...\rbrack\ suis partibus %
\textit{(aa)}~uniforme, et in %
\textit{(bb)}~materiae uniformis %
\bars{(ut duo aeres pedales \textit{erg.}}; et %
\textit{(aaa)}~ambo facta sint et %
\textit{(bbb)}~\textlangle ambo\textrangle\ %
\textit{(ccc)}~ambo facta %
\lbrack...\rbrack\ quadam materia %
\textit{(aaaa)}~hom %
\textit{(bbbb)}~uniformi prius %
\bars{(ut aere communi) \textit{erg.}} sed cujus %
\textit{(aaaaa)}~mater %
\textit{(bbbbb)}~copia red %
\textit{(ccccc)}~quantitas %
\textit{(ddddd)}~qu %
\textit{(eeeee)}~quantitates diversae intra duo mobilia %
\bars{quae constituerunt \textit{erg.~u.~gestr.}} sint redactae, %
\lbrack...\rbrack\ aer pedalis %
\textit{(aaaaa-a)}~duplo plus aeris communis %
\href{\Scan{LH_35_14_02_001v}}{\lbrack1\,v\textsuperscript{o}\rbrack} %
\textit{(bbbbb-b)}~duplo plus %
\textit{(aaaaa-aa)}~aer\textlangle-- --\textrangle\  %
\textit{(bbbbb-bb)}~\textlangle--\textrangle\ %
\textit{(ccccc-cc)}~aeris communis %
\bars{sit \textit{erg.\;Hrsg.}} quam alter) %
\textit{(aaaaa-aaa)}~tunc %
\textit{(aaaaa-aaaa)}~erit sane mobile %
\textit{(bbbbb-bbbb)}~quae %
\textit{(aaaaa-aaaaa)}~sunt rationes quantitatum %
\textit{(bbbbb-bbbbb)}~quantitates materiae prius uniformis intra diversa mobilia (aeris communis) sunt %
\textit{(bbbbb-bbb)}~tunc %
\textit{(aaaaa-aaaa)}~densitates sunt %
\textit{(bbbbb-bbbb)}~densitas est %
\textit{(aaaaa-aaaaa)}~id ipsum in mobilibus, quod %
\textit{(aaaaa-aaaaa-a)}~est in %
\textit{(bbbbb-bbbbb-b)}~sunt %
\textit{(bbbbb-bbbbb)}~ea ipsa in %
\textit{(aaaaa-aaaaa-a)}~mobi %
\textit{(bbbbb-bbbbb-b)}~ipsis mobilibus %
\bars{aut aliis conformibus \textit{erg.}} quae sunt %
\bars{ut \textit{streicht Hrsg.}} materiae constituentis quantitates %
\textit{(aaaaa-aaaaa-aa)}~seu %
\textit{(bbbbb-bbbbb-bb)}~(seu unus aer %
\textit{(aaaaa-aaaaa-aaa)}~communis de %
\textit{(bbbbb-bbbbb-bbb)}~pedalis altero dicitur duplo densior.%
~\textit{L}%
}}
Et si detur aer cujus
\lbrack\textit{Text bricht ab.}\rbrack
\pend%
\vspace{0.75em}%
\newpage%
\begin{edfigure}[0.00em][0.0em]%
\includegraphics[width=0.50\textwidth]%
{gesamttex/edit_VIII-4/images/LH_35_14_02_001_d_001v.pdf}%
\edcaption[1.0em]{%
\lbrack\textit{Fig.\:1, quer und z.T. überschrieben}\rbrack%
\edtext{}{%
\lemma{\lbrack\textit{Fig.\:1}\rbrack}%
\Cfootnote[nonum]{%
Ein in \textit{L} (Bl.\;\href{\Scan{LH_35_14_02_001v}}{1\,v\textsuperscript{o}}) überlieferter, gestr. Entwurf zum Diagramm wird hier nicht wiedergegeben.%
}}
}%
\end{edfigure}%
\vspace{4.25em}%
\pstart%
\noindent%
\centering%
*     *     *
\pend%
\vspace{3.75em}%
\selectlanguage{german}%
\begin{ledgroup}%
\footnotesize%
\pstart%
\noindent%
\textls{Ausführliche Datierungsgründe:}%
\pagelabel{58216_dat}
Der Träger der undatierten Notiz N.\;\ref{58216} weist keine Spuren eines Wasserzeichens auf.
Die Datierung des Textes ergibt sich aus folgenden Beobachtungen:
\pend%
\pstart%
(1)
Das Zettelformat des Trägers deutet auf eine Entstehung der Notiz N.\;\ref{58216} nach Leibnizens Umzug von Paris\indexO{Paris} nach Hannover:\indexO{Hannover}
Kleine und fragmentarische Zettel sind im Nachlass seltener aus der Mainzer\indexO{Mainz} oder Pariser\indexO{Paris} Zeit überliefert.
Folglich ist unwahrscheinlich, dass die Notiz vor Ende Dezember 1676 verfasst wurde.
\pend%
\pstart%
(2)
Die Notiz N.\;\ref{58216} drückt eine intuitive Feststellung aus:
Die Dichte einer gleichförmigen Materie (wie etwa Luft) sei bei gleichem Volumen (und gleicher Temperatur) proportional zur Quantität der Materie.
Dass diese Feststellung in Texten, die spätestens Mitte 1686 entstanden sind (wie etwa N.\;\ref{60212}, \ref{41212} und \ref{41213} im vorliegenden Band), als selbstverständlich vorausgesetzt wird, muss freilich nicht bedeuten, dass N.\;\ref{58216} davor entstand.
Denn die Notiz N.\;\ref{58216} beabsichtigt wohl nicht, einen neue Inhalte darzulegen, sondern vielmehr, einen gemeinverständlichen Gedanken rigorös und anschaulich auszuformulieren.
\pend%
\pstart%
(3)
In dieser Hinsicht ist zu beachten, dass ein ähnlicher Gedanke %
noch in der \textit{Dynamica}, pars\;I, sectio\;I, cap.\;II, def.\;2\cite{01354} anzutreffen ist (\textit{CFP}~II, S.\;36.1--26~[C]). %
Auch dort versucht Leibniz, den intuitiven Begriff der Dichte rigorös und anschaulich auszuformulieren.
Die Textgenese der ältesten Fassung der \textit{Dynamica} zeigt zudem, dass Leibniz mit der Formulierung dieses Begriffs gehadert hat (vgl. \textit{CFP}~I, S.\;28.8--30.26~[A] samt Variantenapparat).
Demnach erweist es sich als plausibel, dass die Notiz N.\;\ref{58216} im Zusammenhang mit der Bearbeitung dieses Teils der \textit{Dynamica}\cite{01354} verfasst wurde, d.\,h. zwischen Sommer und November 1689.
Die italienische Herkunft des Trägers von N.\;\ref{58216} lässt sich aufgrund des kleinen Papierformats weder bestätigen noch zurückweisen;
zahlreiche Manuskripte der Signatur LH~XXXV~14,\;2 sind aber tatsächlich auf italienischem Papier verfasst.
Schließlich sei bemerkt, dass auch die im vorliegenden Band edierten Texte N.\;\ref{58656-A} und N.\;\ref{58258} im Zusammenhang mit der Bearbeitung der ersten Fassung der \textit{Dynamica}\cite{01354} entstanden und mit dem Begriff der Dichte befasst sind.
\pend%
\pstart%
(4)
Eine frühere Datierung der Notiz N.\;\ref{58216} als \textit{Sommer bis November 1689} bleibt aber möglich.%
\pend%
\end{ledgroup}%
\normalsize%
\selectlanguage{latin}%
\frenchspacing%
